\documentclass[11pt]{article}

\usepackage[margin=1.1in]{geometry}
\usepackage{amsmath,amssymb,amsthm,mathtools}
\usepackage{enumitem}
\usepackage{hyperref}
\usepackage[nameinlink,noabbrev]{cleveref}

\hypersetup{colorlinks=true,linkcolor=blue,citecolor=blue,urlcolor=blue}

\newtheorem{theorem}{Theorem}[section]
\newtheorem{lemma}[theorem]{Lemma}
\newtheorem{corollary}[theorem]{Corollary}
\newtheorem{conjecture}[theorem]{Conjecture}
\theoremstyle{remark}
\newtheorem{remark}[theorem]{Remark}

\providecommand{\doi}[1]{\href{https://doi.org/#1}{doi:#1}}

\title{Greedy weakly complete sequences with initial terms $n$ and $n+d$}
\author{Yitzchak Shmalo\\
Einstein Institute of Mathematics, The Hebrew University of Jerusalem\\
Givat Ram, Jerusalem, Israel}
\date{}

\begin{document}
\maketitle

\begin{abstract}
Fix integers $1\le d<n$ and let $a_1=n$, $a_2=n+d$; for $k\ge 2$ let $a_{k+1}$ be the least integer greater than
$a_k$ that is not a sum of distinct terms among $a_1,\dots,a_k$. Knapp and Paul conjectured that after $n+1$ initial
terms this sequence consists of blocks of $n-1$ terms, each a translate of $\{0,1,\dots,n-1\}\setminus\{d-1\}$, with
translations growing geometrically with ratio $n$; they proved the case $d=1$, and Fox and Knapp proved the cases
$2\le d\le n-8$. We prove the remaining cases $n-7\le d\le n-1$, which completes the proof of the conjecture. The
proof describes the set of subset sums exactly at the end of every block and reduces the induction step to an
inclusion between sets of sums of distinct elements of $\{0,\dots,n-1\}\setminus\{d-1\}$, proved directly for
$n\ge 20$ and checked by computer for $n\le 19$.
\end{abstract}

\noindent\textbf{Mathematics Subject Classification (2020):} 11B13, 11B83.

\noindent\textbf{Keywords:} weakly complete sequence, greedy algorithm, subset sums.

\section{Introduction}

An increasing sequence of positive integers is \emph{weakly complete} if every sufficiently large integer is a sum of
distinct terms of the sequence \cite{Hon,KP}. Given two initial terms, the greedy algorithm produces such a sequence:
once $a_1<\dots<a_k$ are known, $a_{k+1}$ is the least integer greater than $a_k$ that is not a sum of distinct terms
among $a_1,\dots,a_k$. Every integer greater than $a_2$ is then a sum of distinct terms. For example, the initial
terms $5,8$ give
\[
5,\ 8,\ 9,\ 10,\ 11,\ 12,\ 48,\ 49,\ 51,\ 52,\ 248,\ 249,\ 251,\ 252,\ 1248,\ 1249,\ 1251,\ 1252,\ \dots
\]
(see \cite{FK}). Knapp and Paul \cite{KP} determined these sequences for several families of initial terms and
observed that the initial terms are best written as $n$ and $n+d$. For $d<n$ they made the following conjecture
\cite[Conjecture~B]{KP}.

\begin{conjecture}[Knapp and Paul]\label{conj:KP}
Let $1\le d<n$ and let $(a_k)$ be the greedy sequence with $a_1=n$ and $a_2=n+d$. Then the set of its terms is
\[
I\cup J\cup\bigcup_{i\ge 1}\bigl(an^i+S\bigr),
\]
where $I=\{n,n+d\}$, $J=\{n+d+1,\dots,2n+d-1\}$,
\[
a=\frac{3n^2+(2d-3)n-4d+2}{2(n-1)},\qquad c=\frac{n^2-n-2(d-1)}{2(n-1)},
\]
and $S=\{-c+j:\ 0\le j\le n-1,\ j\ne d-1\}$.
\end{conjecture}

The case $d=1$ is \cite[Theorem~2]{KP}, and Fox and Knapp proved the conjecture for $2\le d\le n-8$
\cite[Theorem~1]{FK}. For $n=3$ and $d=2$ the sequence is entry A072134 of the OEIS \cite{OEIS}, where the same
description is given. Linz and Jones \cite{LJ} observe that in these greedy sequences no term is a sum of distinct
smaller terms. We prove the remaining cases.

\begin{theorem}\label{thm:main}
\Cref{conj:KP} holds for all $n\ge 2$ and $\max(1,n-7)\le d\le n-1$. Together with \cite[Theorem~1]{FK} and
\cite[Theorem~2]{KP}, it holds for all $1\le d<n$.
\end{theorem}

It is convenient to state the conclusion in a different form. Put
\[
C=n+d,\qquad W=\{0,1,\dots,n-1\}\setminus\{d-1\},\qquad K=|W|=n-1,
\]
and define
\begin{equation}\label{eq:M}
M_1=n+nC+\binom n2-C+1,\qquad M_{i+1}=nM_i+\binom n2-(d-1)\quad(i\ge 1).
\end{equation}
We show that the sequence is
\begin{equation}\label{eq:seq}
n,\ C,\ C+1,\ \dots,\ C+n-1,\ \text{followed by the blocks } M_1+W,\ M_2+W,\ \dots,
\end{equation}
each block listed in increasing order. This is the conjecture. Indeed $\sum_{w\in W}w=\binom n2-(d-1)=(n-1)c$, so
$M_{i+1}+c=n(M_i+c)$, and therefore $M_i=(M_1+c)n^{i-1}-c$. A direct expansion gives
$(M_1+c)/n=\tfrac32 n+d-\tfrac{d-1}{n-1}$, which equals $a$. Hence $M_i+W=an^i+S$, while $\{n\}\cup\{C,\dots,C+n-1\}=I\cup J$.

Knapp and Paul \cite{KP} and Fox and Knapp \cite{FK} argue by induction on the blocks, and so do we. Fox and Knapp
cover the integers below the next block by intervals, among them intervals $D_m$, each made of the sums of $m$ terms
of the current block and one integer that is already a sum of earlier terms. For $2\le d\le n-8$ they show that $D_m$ and
$D_{m+1}$ are contiguous for $2\le m\le n-4$, and they represent the integers between the last two of these intervals
by adding the sum of all earlier terms to sums of terms of the current block \cite[Section~3]{FK}. At the end of that
section they note that for $n<d+5$ gaps already appear between the sets of sums of the initial terms $C,\dots,C+n-1$, so that
the integers in these gaps need explicit representations, and that they expect the formula to hold for all $d<n$
with a more refined analysis. We give these representations for every block and every number of new terms at once.
For a finite set $A$ of positive integers let $\Sigma(A)$ be the set of sums of its subsets, the empty sum $0$
included, and let $T$ be the sum of the elements of $A$. We show that at the end of every block $\Sigma(A)=R_T$,
where for an integer $T\ge 2C$
\begin{equation}\label{eq:RT}
R_T=\{0,n\}\cup[C,T-C]\cup\{T-n,T\}
\end{equation}
and $[x,y]$ denotes the set of integers $m$ with $x\le m\le y$ (\Cref{lem:start,lem:step}). The induction step then
reduces to the inclusion \eqref{eq:gap} between sets of sums of distinct elements of $W$, which involves neither the
block index nor the terms already chosen. For $n\ge 20$ and $d\ge n-7$ the inclusion is proved in
\Cref{lem:analytic}, and for $n\le 19$ it is checked by computer.

\section{The subset sums at the end of a block}

Let $\Sigma(A)$, $T$, $[x,y]$ and $R_T$ be as in \eqref{eq:RT} and the sentences around it. The set $\Sigma(A)$ is
symmetric under $x\mapsto T-x$.

\begin{lemma}\label{lem:start}
The greedy sequence begins with $n,C,C+1,\dots,C+n-1$. If $A_0$ is the set of these $n+1$ terms and
$T_0=n+nC+\binom n2$ is their sum, then $\Sigma(A_0)=R_{T_0}$.
\end{lemma}

\begin{proof}
Every sum of two or more distinct terms that includes a term at least $C$ is at least $C+n$. So for $C\le x\le
C+n-1$ the only sums of earlier terms that are at most $x$ are $0$, $n$ and the single terms, and each such $x$ is
chosen in turn.

A sum using $r$ of the terms $C,\dots,C+n-1$ has the form $rC+q$, where $q$ is a sum of $r$ distinct elements of
$\{0,\dots,n-1\}$. These $q$ fill the interval $[\binom r2,\binom r2+r(n-r)]$: from any $r$-subset other than the
largest, some element can be raised by one without leaving $\{0,\dots,n-1\}$ or meeting another element. Write
$l_r=rC+\binom r2$ and $u_r=l_r+r(n-r)$. Adding the optional term $n$, the sums with exactly $r$ terms from the interval
form $[l_r,u_r]\cup[l_r+n,u_r+n]$. For $1\le r\le n-1$ we have $u_r-l_r\ge n-1$, so this union is the interval
$[l_r,u_r+n]$. Moreover
\[
l_{r+1}-(u_r+n)-1=d-1-r(n-1-r)\le 0\qquad(1\le r\le n-2),
\]
because $r(n-1-r)\ge n-2\ge d-1$. So the bands for $1\le r\le n-1$ together form the interval $[l_1,u_{n-1}+n]$. Here
$l_1=C$, and $u_{n-1}+n=T_0-C$ by a direct computation. The case $r=0$ contributes $\{0,n\}$, and $r=n$ contributes
$\{T_0-n,T_0\}$.
\end{proof}

For $0\le r\le K$ let $Q_r$ be the set of sums of $r$ distinct elements of $W$, with $Q_0=\{0\}$, and let
$l(r)=\min Q_r$ and $h(r)=\max Q_r$.

\begin{lemma}\label{lem:Q}
Any two consecutive elements of $Q_r$ differ by at most $2$.
\end{lemma}

\begin{proof}
Let $B\subset W$ have $r$ elements and not be the set of the $r$ largest elements of $W$. Then some $b\in B$ has a
successor $b'$ in $W$ that is not in $B$, and replacing $b$ by $b'$ raises the sum by $b'-b\in\{1,2\}$, since $W$
omits only one integer of $[0,n-1]$. Starting from the $r$ smallest elements and repeating this, we reach the
$r$ largest. The sums along the way go from $l(r)$ to $h(r)$ in steps of $1$ or $2$.
\end{proof}

For $d\ge 2$ a precise description of the sets of sums of $r$ distinct elements of $W$ is given by Fox and Knapp
\cite[Section~2]{FK}.

The induction step rests on the following inclusion. For $0\le r<K$ consider
\begin{equation}\label{eq:gap}
[h(r),\ l(r+1)+C-1]\ \subset\ \bigl(Q_r+\{d-1,\,C-1\}\bigr)\cup\bigl(Q_{r+1}+\{0,n\}\bigr).
\end{equation}
When $h(r)>l(r+1)+C-1$ the left side is empty and there is nothing to check.

\begin{lemma}\label{lem:step}
Suppose \eqref{eq:gap} holds for every $0\le r<K$. Let $A$ be the set of the first terms of the greedy sequence up to
some point, let $T$ be the sum of these terms, and suppose that $\Sigma(A)=R_T$, $T\ge 2C+1$ and that the largest element of $A$
is less than $M=T-C+1$. Then the next $K$ terms of the sequence are the elements of $M+W$ in increasing order. The new
sum of the terms is $T'=T+KM+\sum_{w\in W}w$, and $\Sigma(A\cup(M+W))=R_{T'}$. Moreover $T'-C+1=nM+\binom n2-(d-1)$, and this
number exceeds $M+n-1$.
\end{lemma}

\begin{proof}
Since $T-C=M-1$, $T-n=M+d-1$ and $T=M+C-1\ge M+n$, the only element of $R_T$ in $[M,M+n-1]$ is $M+d-1$, and every
integer in $[C,M-1]$ lies in $R_T$. Let $x$ be a new term in $[M,M+n-1]$. Then $x$ plus a positive element of $R_T$ is
at least $M+n$, and the sum of two new terms is at least $2M>M+n-1$. So no new term creates a sum in
$[M,M+n-1]$, and the greedy algorithm chooses every integer of $[M,M+n-1]$ other than $M+d-1$. These are the elements
of $M+W$.

Next we show that $\Sigma(A')=R_{T'}$, where $A'=A\cup(M+W)$. A subset of $A'$ that uses $r$ of the new terms has sum
$rM+q+s$ with $q\in Q_r$ and $s\in R_T$. Hence
\[
\Sigma(A')=\bigcup_{r=0}^{K}\bigl(rM+Q_r+R_T\bigr).
\]
The interval $[C,T-C]$ contains at least two integers, so by \Cref{lem:Q} the set $Q_r+[C,T-C]$ is the interval
$[l(r)+C,\,h(r)+T-C]$. Thus the sums with exactly $r$ new terms and a middle element of $R_T$ form the interval
\[
D_r=[rM+l(r)+C,\ (r+1)M+h(r)-1].
\]
These are the intervals $D_m$ of \cite[Section~3]{FK}. The remaining sums with $r$ new terms are $rM+Q_r+\{0,n\}$ and
$rM+Q_r+\{T-n,T\}=(r+1)M+Q_r+\{d-1,C-1\}$. The integers strictly between $D_r$ and $D_{r+1}$ form the set
$(r+1)M+[h(r),l(r+1)+C-1]$. By \eqref{eq:gap} this set is contained in
$(r+1)M+(Q_r+\{d-1,C-1\})\cup(Q_{r+1}+\{0,n\})$, which lies in $\Sigma(A')$. For $r=K-1$ this is the
representation used in \cite[Section~3]{FK}. Hence $\Sigma(A')$ contains the
interval from $\min D_0=C$ to $\max D_K=(K+1)M+h(K)-1=T'-C$. Every element of $\Sigma(A')$ other than those in
$R_T$ is at least $M>C$, so the only elements of $\Sigma(A')$ below $C$ are $0$ and $n$. By the symmetry
$x\mapsto T'-x$ of $\Sigma(A')$, the only elements above $T'-C$ are $T'-n$ and $T'$. Therefore
$\Sigma(A')=R_{T'}$.

Finally, $T'-C+1=M+KM+\sum_{w\in W}w=nM+\binom n2-(d-1)$, which exceeds $M+n-1$ because $M\ge n$.
\end{proof}

\begin{proof}[Proof of \eqref{eq:seq}, assuming \eqref{eq:gap}]
After the first $n+1$ terms, \Cref{lem:start} shows that the hypotheses of \Cref{lem:step} hold with
$T=T_0$ and $M=T_0-C+1=M_1$; here $T_0-2C\ge 1$ and $M_1-(C+n-1)=(n-2)C+\binom n2+2>0$. \Cref{lem:step} gives the block $M_1+W$
and again a state $\Sigma(A)=R_T$ with $T$ larger and with the next value of $M$ equal to $M_2$, which exceeds the last
term. After the block, every integer from $M+n$ to the next $M$ minus one lies in $R_{T'}$, so no term is chosen
there, and the block beginning at the next $M$ follows. Induction on $i$ gives \eqref{eq:seq}, with the recurrence
\eqref{eq:M}.
\end{proof}

\section{Proof of the inclusion}

Write $\sigma=\sum_{w\in W}w$. Taking complements $B\mapsto W\setminus B$ shows that $Q_{K-j}=\sigma-Q_j$ for
$0\le j\le K$, so that $h(K-1-r)=\sigma-l(r+1)$ and $l(K-r)=\sigma-h(r)$. Let $r'=K-1-r$ and consider the reflection
$\rho(x)=\sigma+C-1-x$. It maps the left side $[h(r),\,l(r+1)+C-1]$ of \eqref{eq:gap} onto
$[h(r'),\,l(r'+1)+C-1]$, the left side of \eqref{eq:gap} for $r'$. Using $C-d=n$, it maps the four translates on the
right side as follows:
\begin{gather*}
\rho(Q_r+d-1)=Q_{r'+1}+n,\qquad \rho(Q_r+C-1)=Q_{r'+1},\\
\rho(Q_{r+1})=Q_{r'}+C-1,\qquad \rho(Q_{r+1}+n)=Q_{r'}+d-1.
\end{gather*}
Hence \eqref{eq:gap} for $r$ and for $K-1-r$ are equivalent, and it is enough to prove \eqref{eq:gap} for
$r\le (K-1)/2$, that is, for $s=r+1\le n/2$.

\begin{lemma}\label{lem:analytic}
If $n\ge 20$ and $n-7\le d\le n-1$, then \eqref{eq:gap} holds for $0\le r\le (K-1)/2$.
\end{lemma}

\begin{proof}
For $r=0$ the left side of \eqref{eq:gap} is $[0,C-1]$, since $l(1)=0$. The set $W\cup(n+W)$ is $[0,2n-1]$ with
$d-1$ and $n+d-1=C-1$ removed, and these two numbers form $Q_0+\{d-1,C-1\}$.

Let $2\le s=r+1\le n/2$. Since $d\ge n-7$, the set $W$ contains $[0,d-2]$, which has $d-1\ge n-8$ elements, and
$n-8\ge n/2\ge s$ because $n\ge 16$.
Sums of $s$ distinct elements of $[0,d-2]$ fill the interval $[\binom s2,\binom s2+s(d-1-s)]$, by the argument in the
proof of \Cref{lem:start}. In particular $l(s)=\binom s2$. The function $s\mapsto s(n-8-s)$ is concave, and at
$s=2$ and $s=n/2$ it equals $2n-20$ and $\tfrac n2(\tfrac n2-8)$, both at least $n$ when $n\ge 20$. Hence
$s(d-1-s)\ge n$, and $Q_s\cup(n+Q_s)$ contains $[l(s),l(s)+2n]$, which contains $[l(s),l(s)+C-1]$ because $C-1\le
2n-2$. Finally $h(s-1)\ge 1+2+\dots+(s-1)=l(s)$, since $W$ contains $1,\dots,s-1$. So the left side of
\eqref{eq:gap} lies in $Q_{r+1}+\{0,n\}$.
\end{proof}

For $2\le n\le 19$ we check \eqref{eq:gap} by computer for every $1\le d<n$ and every $0\le r<K$. For each $(n,d)$
the sets $Q_0,\dots,Q_K$ are computed by the update $Q_j\leftarrow Q_j\cup(w+Q_{j-1})$, applied for each $w\in W$
with $j$ decreasing, starting from $Q_0=\{0\}$ and $Q_j=\varnothing$ for $j>0$. Then every integer between $h(r)$ and
$l(r+1)+C-1$ is tested against the four translates in \eqref{eq:gap}. Each of the $171$ pairs $(n,d)$ passes. This is
a finite computation that covers the whole induction, since \eqref{eq:gap} depends only on $n$, $d$ and $r$.

\begin{proof}[Proof of \Cref{thm:main}]
For $n\ge 20$ and $n-7\le d\le n-1$, \Cref{lem:analytic} and the reflection give \eqref{eq:gap} for all $0\le r<K$.
For $n\le 19$ the computation gives it. In both cases \eqref{eq:seq} holds, which is the conjecture. The remaining
pairs are covered by \cite[Theorem~2]{KP} for $d=1$ and \cite[Theorem~1]{FK} for $2\le d\le n-8$.
\end{proof}

\begin{remark}
The argument of \Cref{lem:step} works for every $1\le d<n$; only the verification of \eqref{eq:gap} uses the
restriction on $d$.
\end{remark}

\section*{Acknowledgments and funding}

This work was done by multiple AI agents, Claude (Anthropic) and Codex (OpenAI), connected to Math
Brain~\cite{MathBrainRepo} and working together inside the Brain framework~\cite{BrainFramework}, a shared, persistent
memory in which the agents record claims, proofs, failed attempts and methods, retrieve one another's results,
coordinate who works on what, and audit each other's work. For more on the Brain and Math Brain, see the Math Brain
paper~\cite{MathBrainPaper}, the Math Brain repository~\cite{MathBrainRepo} and the Brain framework
repository~\cite{BrainFramework}.

The research presented in this paper was supported by the European Research Council (ERC) under the European
Union's Horizon Europe research and innovation programme (grant agreement No.\ 101041711). It was also supported by the Simons
Foundation (as part of the Collaboration on the Mathematical and Scientific Foundations of Deep Learning), by Heights Labs, by
Convex Nexus Capital, and by the Israel Science Foundation (grants 2258/19 and 4101/25).

\section*{Declaration of competing interest}

The author holds equity in Heights Labs and Convex Nexus Capital and has a management role at Convex Nexus Capital.

\section*{Declaration of generative AI and AI-assisted technologies}

The reduction to \eqref{eq:gap} and its verification were developed with AI systems (OpenAI Codex and Anthropic Claude
agents working under the author's direction), which also drafted the text and wrote the computer checks. The
conjecture and its formula are due to Knapp and Paul \cite{KP}, who proved the case $d=1$ \cite[Theorem~2]{KP}, and
the cases $2\le d\le n-8$ to Fox and Knapp \cite{FK}, whose intervals $D_m$ and representation of the integers below
a block are used in \Cref{lem:step}. The author directed the work and is responsible for the content of the paper and for decisions concerning its
dissemination.

\section*{Data and code availability}

The Python program that verifies \eqref{eq:gap} for $2\le n\le 19$ is available at
\url{https://github.com/yspennstate/greedy-weakly-complete-sequences/tree/v1.0}.

\begin{thebibliography}{9}

\bibitem{FK}
A.~Fox and M.~P.~Knapp,
\emph{A note on weakly complete sequences},
J.\ Integer Seq.\ \textbf{16} (2013), Article 13.6.2.
\url{https://cs.uwaterloo.ca/journals/JIS/VOL16/Knapp/knapp5.html}

\bibitem{Hon}
R.~Honsberger,
\emph{Mathematical Gems III},
Mathematical Association of America, 1985.

\bibitem{KP}
M.~P.~Knapp and M.~D.~Paul,
\emph{On weakly complete sequences formed by the greedy algorithm},
Integers \textbf{12} (2012), A4.
\url{https://math.colgate.edu/~integers/m4/m4.pdf}

\bibitem{LJ}
W.~Linz and E.~Jones,
\emph{$r$-completeness of sequences of positive integers},
Integers \textbf{16} (2016), A59.

\bibitem{OEIS}
OEIS Foundation Inc.,
\emph{The On-Line Encyclopedia of Integer Sequences}, sequence A072134, accessed November 2026.
\url{https://oeis.org/A072134}

\bibitem{MathBrainRepo}
Y.~Shmalo, \emph{Math Brain: public snapshot, code and tests}, 2026.
\url{https://github.com/yspennstate/math-brain-public}

\bibitem{MathBrainPaper}
Y.~Shmalo, \emph{Math Brain: shared, auditable memory for multi-agent mathematical research}, October 2026.
\url{https://github.com/yspennstate/math-brain-public/tree/main/papers/math-brain-2026-10-24}

\bibitem{BrainFramework}
Y.~Shmalo, \emph{Brain Framework}, public software repository, 2026.
\url{https://github.com/yspennstate/brain-framework-public}

\end{thebibliography}

\end{document}
